% ============================================================
% Conteúdo original do metadata.yaml (migrado para cá)
% ============================================================
\makeatletter
\renewcommand\paragraph{\@startsection{paragraph}{4}{\z@}%
  {-3.25ex\@plus -1ex \@minus -.2ex}%
  {1.5ex \@plus .2ex}%
  {\normalfont\normalsize\bfseries}}
\makeatother
\usepackage{caption}
\captionsetup[table]{labelformat=empty}
\usepackage{newfloat}
\DeclareFloatingEnvironment[fileext=loquad,listname={Índice de Quadros},name=Quadro]{quadro}
\captionsetup[quadro]{labelformat=empty}
\DeclareFloatingEnvironment[fileext=lotab,listname={Índice de Tabelas},name=Tabela]{tabela}
\captionsetup[tabela]{labelformat=empty}
\usepackage{chngcntr}
\counterwithin{figure}{section}
\renewcommand{\figurename}{Figura}
\usepackage{pdflscape}
\newcommand{\blandscape}{\begin{landscape}}
\newcommand{\elandscape}{\end{landscape}}
\usepackage{fancyhdr}
\pagestyle{fancy}
\fancyhf{}
\cfoot{\thepage}
\renewcommand{\headrulewidth}{0pt}
\fancypagestyle{plain}{%
  \fancyhf{}%
  \cfoot{\thepage}%
  \renewcommand{\headrulewidth}{0pt}%
}
\usepackage{listings}
\usepackage{xcolor}
\definecolor{codebg}{rgb}{0.95,0.95,0.95}
\lstset{
  basicstyle=\ttfamily\footnotesize,
  breaklines=true,
  backgroundcolor=\color{codebg},
  frame=single,
  rulecolor=\color{codebg},
  showstringspaces=false
}

% ============================================================
% Melhorias tipográficas para tabelas (novo)
% ============================================================
\usepackage{booktabs}
\usepackage{microtype}
\usepackage{array}
\usepackage{etoolbox}
\usepackage[htt]{hyphenat}

% Espaçamento vertical nas tabelas
\renewcommand{\arraystretch}{1.3}

% Permitir quebra de linha em sublinhados (underscores) no modo texto
\renewcommand{\_}{\textunderscore\allowbreak}

% Reduzir o tamanho da fonte nas tabelas geradas pelo Pandoc (longtable)
% para ~80% do tamanho normal, dando mais espaço ao conteúdo
\AtBeginEnvironment{longtable}{\small}
